% Computer Science A --- Spring Boot (\texttt{application.*}, auto-config, actuator, packaging, Boot testing). See \texttt{spring.tex} for the core framework.
% Free-response bank at the front is migrated from \texttt{cs-a/06-practice.tex} (Spring Boot/configuration, Lombok, Actuator).

\runningheader{Computer Science A}{Spring Boot}{Page \thepage\ of \numpages}

\begingradingrange{csaspringboot}

\uplevel{\par\textbf{Free response: Spring Boot, configuration, Lombok, and Actuator}\par
\textit{These items mirror the unit bank; use the boxed space for student work. The printed key expands each idea for study.}\par}

\uplevel{\par\textbf{Spring Boot and configuration}\par}

\question[4]
What is Spring Boot?
\begin{solutionorbox}[1.25in]
  \textbf{Spring Boot} is an opinionated layer \emph{on top of} the Spring Framework: \textbf{dependency starters} pull coordinated libraries, \textbf{auto-configuration} wires common beans when classpath and properties match, and an \textbf{embedded HTTP server} (Tomcat by default) lets you run \texttt{java -jar} with minimal XML or Java config boilerplate.

  It does \textbf{not} replace Spring concepts (beans, DI, MVC, Data); it accelerates getting to a working application and standardizes production concerns (metrics, health, externalized config, executable JAR packaging).
\end{solutionorbox}

\question[4]
Tell me about the \texttt{@SpringBootApplication} annotation.
\begin{solutionorbox}[1.35in]
  \texttt{@SpringBootApplication} is a \textbf{convenience composed annotation} used on the class that contains \texttt{public static void main}. It bundles:
  \begin{itemize}
    \item \texttt{@Configuration} --- the class can declare \texttt{@Bean} methods.
    \item \texttt{@EnableAutoConfiguration} --- pull in Boot's conditional auto-configuration classes.
    \item \texttt{@ComponentScan} (by default from the package of the annotated class downward) --- discover \texttt{@Component}, \texttt{@Service}, \texttt{@Repository}, \texttt{@Controller}, etc.
  \end{itemize}
  Together, these bootstrap a Boot application's application context the way the course expects you to describe ``the main class.''
\end{solutionorbox}

\question[4]
What is Spring Initializr?
\begin{solutionorbox}[1.1in]
  \textbf{Spring Initializr} (\texttt{start.spring.io}, also integrated into IDEs) generates a \textbf{new project skeleton}: build file (\texttt{pom.xml} or \texttt{build.gradle}), default package layout, main application class, test stub, and chosen \textbf{dependencies} (Web, Data JPA, Validation, Actuator, etc.).

  It encodes sensible defaults (Java version, Spring Boot BOM, plugin configuration) so learners start from a compiling project instead of hand-assembling classpath coordinates.
\end{solutionorbox}

\question[4]
What is the use of a \texttt{.properties} or \texttt{.yaml} file?
\begin{solutionorbox}[1.35in]
  Files such as \texttt{application.properties} or \texttt{application.yml} \textbf{externalize configuration}: ports, context paths, datasource URLs, feature flags, logging levels, and third-party API keys (often \emph{referenced} here but supplied via environment variables in real deployments).

  Boot maps dotted keys (or YAML hierarchies) onto beans via relaxed binding and \texttt{@ConfigurationProperties}. \textbf{Profiles} (\texttt{application-dev.yml}, \texttt{spring.profiles.active}) let you swap whole configuration bundles per environment without recompiling.
\end{solutionorbox}

\newpage

\uplevel{\par\textbf{Lombok and Actuator}\par}

\question[4]
What is Lombok and its benefits?
\begin{solutionorbox}[1.25in]
  \textbf{Project Lombok} is an annotation processor that generates repetitive Java bytecode at compile time: getters/setters, constructors, \texttt{equals}/\texttt{hashCode}, \texttt{toString}, builders (\texttt{@Builder}), and logging fields (\texttt{@Slf4j}).

  \noindent\textbf{Benefits:} less boilerplate in entities and DTOs, fewer merge conflicts in giant hand-written methods, and readable domain classes.

  \noindent\textbf{Caveat:} IDEs and CI must enable annotation processing so navigation and refactoring stay accurate; some teams limit Lombok features (e.g.\ avoid \texttt{@Data} on JPA entities) to reduce surprising behavior.
\end{solutionorbox}

\question[4]
What is the Spring Actuator module?
\begin{solutionorbox}[1.2in]
  \textbf{Spring Boot Actuator} adds \textbf{production-ready operational endpoints}: health checks for orchestrators, application \texttt{info}, metrics for Prometheus/Micrometer, environment and configuration introspection, thread dumps, loggers control, and more---depending on what you include and expose.

  It complements your business REST API: think \emph{how is the process behaving in production?} rather than \emph{what does a customer transaction do?}
\end{solutionorbox}

\question[4]
What are actuator endpoints?
\begin{solutionorbox}[1.1in]
  An \textbf{actuator endpoint} is a named operational capability (each has an \texttt{id} like \texttt{health}, \texttt{metrics}, \texttt{env}) implemented as a bean behind Spring's management infrastructure.

  Endpoints can be reached over \textbf{HTTP} (commonly under \texttt{/actuator/\ldots} with Boot defaults), \textbf{JMX}, or other technologies if configured. They should be \textbf{secured} in production because they reveal internals.
\end{solutionorbox}

\question[4]
What does enabling an endpoint mean?
\begin{solutionorbox}[1.05in]
  \textbf{Enabling} means the endpoint's implementation is \textbf{active in the application context}: its operations exist and can execute when invoked, subject to security.

  Boot controls this with properties such as \texttt{management.endpoint.<id>.enabled} (defaults differ by endpoint and version). Disabled endpoints are absent or inert even if their code is on the classpath.
\end{solutionorbox}

\question[4]
What does exposing an endpoint mean?
\begin{solutionorbox}[1.05in]
  \textbf{Exposing} selects \emph{which transport} may reach an enabled endpoint---for example, including an endpoint id in \texttt{management.endpoints.web.exposure.include} so it appears on the HTTP management port/prefix.

  An endpoint can be \textbf{enabled but not exposed} over HTTP (still reachable via JMX if exposed there), which is the modern default posture for sensitive introspection endpoints.
\end{solutionorbox}

\question[4]
What does it mean when an endpoint is ``available''?
\begin{solutionorbox}[1.1in]
  In Boot documentation, \textbf{available} colloquially means the endpoint is both \textbf{enabled} and \textbf{exposed} on the interface you are using (HTTP or JMX), so a client \emph{can} call it assuming security permits.

  If an endpoint is disabled or not included for web exposure, clients receive \texttt{404} on the web adapter even though related beans might exist elsewhere.
\end{solutionorbox}

\question[4]
What are some endpoints that actuator provides and their use?
\begin{solutionorbox}[1.45in]
  Representative endpoints (exact availability depends on classpath and configuration):
  \begin{itemize}
    \item \texttt{/actuator/health} --- liveness/readiness for load balancers and Kubernetes probes.
    \item \texttt{/actuator/info} --- static build or git metadata for dashboards.
    \item \texttt{/actuator/metrics} / Micrometer registry --- throughput, latency, JVM, and custom timers.
    \item \texttt{/actuator/env} and \texttt{/actuator/configprops} --- debug configuration drift (high sensitivity).
    \item \texttt{/actuator/loggers} --- adjust log levels at runtime (use with care).
  \end{itemize}
  Treat anything beyond \texttt{health} (and maybe \texttt{info}) as \textbf{privileged operations} in production.
\end{solutionorbox}

\question[4]
How can you create a custom actuator endpoint?
\begin{solutionorbox}[1.25in]
  Define a \texttt{@Endpoint} bean (Boot 2+ style) with a unique id, expose read/write operations with \texttt{@ReadOperation}, \texttt{@WriteOperation}, or \texttt{@DeleteOperation}, and register it so the actuator infrastructure can adapt it to HTTP/JMX.

  \noindent\textbf{Contrast:} a normal \texttt{@RestController} serves \emph{application} APIs; actuator endpoints follow the management model so they participate in global enable/expose/security policies. For Boot 3/Jakarta, use the \texttt{org.springframework.boot.actuate.endpoint.annotation} package; secure custom endpoints like you would \texttt{/actuator/env}.
\end{solutionorbox}

\newpage

\question[4]
Overall project configuration can be provided to a Spring application using which files?
\begin{choices}
  \choice \texttt{application.xml} and \texttt{application.yml}
  \choice \texttt{application.properties} and \texttt{application.xml}
  \choice \texttt{application.xml} and \texttt{bootstrap.xml}
  \CorrectChoice \texttt{application.yml} and \texttt{application.properties}
\end{choices}
\begin{solution}
  Spring Boot convention loads \textbf{\texttt{application.properties}} or \textbf{\texttt{application.yml}} (and profile-specific variants) from the classpath. \texttt{bootstrap} files belong to Spring Cloud Context, not the baseline pair in the stem.
\end{solution}

\question[4]
Using the default configuration, Spring Boot includes an embedded
\begin{choices}
  \CorrectChoice Tomcat server
  \choice Postgres server
  \choice in-memory database
  \choice user
\end{choices}
\begin{solution}
  The default embedded servlet container is \textbf{Apache Tomcat}. You can switch to Jetty or Undertow via starters. An in-memory database is optional via dependencies (e.g.\ H2), not the default embedded HTTP stack.
\end{solution}

\question[4]
Spring Boot features an ``opinionated configuration'' which means that:
\begin{choices}
  \choice It requires you to tell it how to configure all functionality.
  \CorrectChoice When not provided with configuration, it uses reasonable default configurations.
  \choice It cannot be configured.
  \choice You will never need to configure it.
\end{choices}
\begin{solution}
  \textbf{Opinionated defaults} (starters, auto-configuration) get a project running quickly; you override only what you need. You can still configure extensively---it is not ``unconfigurable.''
\end{solution}

\question[4]
Spring Boot actuators can be used to
\begin{choices}
  \choice interact with your server using a convenient GUI.
  \choice open a running application on the command line.
  \CorrectChoice get runtime information about the application via HTTP.
  \choice start a server when it is not running.
\end{choices}
\begin{solution}
  \textbf{Actuators} expose operational endpoints (health, metrics, env, \ldots), often over HTTP (and JMX), for monitoring and management. They are not a general GUI, a shell, or a server launcher.
\end{solution}

\question[4]
The Spring Boot actuator automatically restarts your Spring application whenever files on the classpath change.
\begin{choices}
  \choice TRUE
  \CorrectChoice FALSE
\end{choices}
\begin{solution}
  Automatic restart on classpath changes is provided by \textbf{Spring Boot DevTools}, not by Actuator alone. Actuator focuses on observability and management endpoints.
\end{solution}

\question[4]
By default all the actuator endpoints are enabled except for \texttt{shutdown}.
\begin{choices}
  \CorrectChoice TRUE
  \choice FALSE
\end{choices}
\begin{solution}
  This follows the classic teaching summary: many endpoints are exposed by default in older defaults; \texttt{shutdown} is sensitive and typically off unless explicitly enabled. Exact exposure changed across Boot 2.x/3.x (e.g.\ only \texttt{health} over HTTP unless configured); treat this as the keyed conceptual answer and align with your Boot version in lecture.
\end{solution}

\question[4]
What kind of file will a Spring Boot project typically be packaged as?
\begin{choices}
  \choice EXE file
  \choice BAT file
  \choice CMD file
  \CorrectChoice JAR file
\end{choices}
\begin{solution}
  Executable \textbf{fat JAR}s (\texttt{spring-boot-maven-plugin} / Gradle plugin) bundle dependencies and \texttt{Main-Class} for \texttt{java -jar}. WAR deployment to an external servlet container is also common but the stem targets the typical standalone artifact.
\end{solution}

\question[4]
What is Spring Boot autoconfiguration?
\begin{choices}
  \choice Autoconfiguration automatically wires up your bean dependencies and adds them to the IOC container
  \CorrectChoice Autoconfiguration attempts to automatically configure your Spring application based on the jar dependencies that you have added
  \choice Autoconfiguration sets up REST API endpoints based on the classes configured in your application
  \choice Autoconfiguration sets the Java version that your Spring project will use based on the features that your code uses
\end{choices}
\begin{solution}
  \textbf{@EnableAutoConfiguration} (via \texttt{@SpringBootApplication}) pulls in conditional configuration classes on the classpath (e.g.\ DataSource when JDBC starter + properties present). DI wiring is broader than auto-config alone; REST endpoints come from your controllers, not auto-config wholesale.
\end{solution}

\question[4]
Spring Boot Actuator custom endpoints can be created by\ldots
\begin{choices}
  \choice Writing MockMVC integration tests with custom data
  \choice Writing a controller that uses the \texttt{@RequestMapping} annotation
  \choice Putting the custom endpoint path in the \texttt{application.yml} file and linking to the implementing class
  \CorrectChoice Annotating a bean with the \texttt{@Endpoint} annotation and give an \texttt{id} property to specify the endpoint path
\end{choices}
\begin{solution}
  Custom \textbf{Actuator} endpoints (JMX or HTTP, depending on exposure settings) use the actuator programming model (\texttt{@Endpoint}, \texttt{@ReadOperation}, etc.). A normal \texttt{@RestController} serves application APIs, not actuator infrastructure.
\end{solution}

\question[4]
How can you test asynchronous code in Spring Boot?
\begin{choices}
  \choice Use the \texttt{@Async} annotation
  \choice Use the \texttt{@TestAsync} annotation
  \CorrectChoice Use the \texttt{CompletableFuture} class
  \choice Asynchronous testing is not supported in Spring Boot
\end{choices}
\begin{solution}
  The keyed answer follows the source. In practice, \texttt{@Async} enables async execution; testing often uses \texttt{CompletableFuture.get()}, \texttt{@Test} timeouts, or Spring's async test utilities (\texttt{WebTestClient}, Awaitility, etc.). There is no standard \texttt{@TestAsync} in core Spring Test.
\end{solution}

\question[4]
How can you customize the test environment using profiles in Spring Boot?
\begin{choices}
  \CorrectChoice Use the \texttt{@ActiveProfiles} annotation
  \choice Use the \texttt{@TestConfiguration} annotation
  \choice Use the \texttt{application-test.properties} file
  \choice Profiles cannot be used to customize the test environment
\end{choices}
\begin{solution}
  \texttt{@ActiveProfiles("test")} on a test class activates profile-specific beans and property files (e.g.\ \texttt{application-test.yml}). \texttt{@TestConfiguration} supplies extra beans; \texttt{application-test.properties} is common but the annotation directly drives which profiles load in tests.
\end{solution}

\question[4]
What is the primary purpose of Spring Boot Testing?
\begin{choices}
  \choice To write production code
  \CorrectChoice To test individual units of code
  \choice To design the application architecture
  \choice To deploy the application in production
\end{choices}
\begin{solution}
  The keyed answer follows the source (unit-focused wording). Spring Boot's test stack also supports \textbf{slice} tests (\texttt{@WebMvcTest}, \texttt{@DataJpaTest}) and full \textbf{integration} tests (\texttt{@SpringBootTest}); clarify in class that ``testing'' spans multiple levels.
\end{solution}

\question[4]
Which testing framework is integrated by default in Spring Boot for writing tests?
\begin{choices}
  \CorrectChoice JUnit
  \choice TestNG
  \choice Mockito
  \choice Cucumber
\end{choices}
\begin{solution}
  \texttt{spring-boot-starter-test} brings in \textbf{JUnit Jupiter} (JUnit 5), AssertJ, Mockito, Hamcrest, \texttt{spring-test}, etc. TestNG and Cucumber are optional additions.
\end{solution}

\question[4]
Explain the role of the \texttt{@WebMvcTest} annotation in Spring Boot Testing.
\begin{choices}
  \CorrectChoice It is used for testing web layers only
  \choice It is used for testing database interactions
  \choice It is used for testing service layers only
  \choice It is used for testing repositories
\end{choices}
\begin{solution}
  \texttt{@WebMvcTest} slices the context to \textbf{web layer} components (controllers, JSON converters, \texttt{MockMvc}) and excludes full data layer auto-configuration unless you import extras. Full-stack tests use \texttt{@SpringBootTest}.
\end{solution}

\question[4]
Explain the purpose of the \texttt{@AutoConfigureMockMvc} annotation in Spring Boot Integration Testing.
\begin{choices}
  \choice It configures the application to use a mock database
  \CorrectChoice It auto-configures the \texttt{MockMvc} instance for testing
  \choice It configures the application to use a real database
  \choice It auto-configures the application context for unit testing
\end{choices}
\begin{solution}
  \texttt{@AutoConfigureMockMvc} registers \texttt{MockMvc} for sending fake HTTP requests in tests (often with \texttt{@SpringBootTest}). It does not imply mock vs.\ real DB; that is controlled by test slices, \texttt{@MockBean}, \texttt{@DataJpaTest}, or Testcontainers.
\end{solution}

\uplevel{\par\textbf{Spring Boot practice bank}\par}

\question[4]
Which annotation usually serves as the primary entry-point annotation for a Spring Boot app?
\begin{choices}
  \CorrectChoice \texttt{@SpringBootApplication}
  \choice \texttt{@EnableJpaRepositories}
  \choice \texttt{@EnableScheduling}
  \choice \texttt{@ConfigurationPropertiesScan} only
\end{choices}
\begin{solution}
  \texttt{@SpringBootApplication} combines \texttt{@Configuration}, \texttt{@EnableAutoConfiguration}, and component scanning, making it the common application bootstrap annotation.
\end{solution}

\question[4]
How does Spring Boot usually decide whether to auto-configure a bean?
\begin{choices}
  \choice It always creates every possible bean from \texttt{spring-boot-autoconfigure}
  \CorrectChoice It uses conditions (classpath, existing beans, properties, environment) to activate matching auto-configuration classes
  \choice It only reads XML files from \texttt{/config}
  \choice It requires manual \texttt{new} calls in the main method
\end{choices}
\begin{solution}
  Auto-configuration is conditional, commonly via \texttt{@ConditionalOnClass}, \texttt{@ConditionalOnMissingBean}, and property-based conditions.
\end{solution}

\question[4]
Which property best sets the HTTP server port to 9090 in Spring Boot?
\begin{choices}
  \CorrectChoice \texttt{server.port=9090}
  \choice \texttt{spring.http.port=9090}
  \choice \texttt{tomcat.server.port=9090}
  \choice \texttt{boot.port=9090}
\end{choices}
\begin{solution}
  \texttt{server.port} is the standard Boot server property in \texttt{application.properties} (or YAML equivalent).
\end{solution}

\question[4]
In current Boot defaults, which Actuator endpoint is typically exposed over HTTP out of the box?
\begin{choices}
  \CorrectChoice \texttt{/actuator/health}
  \choice \texttt{/actuator/shutdown}
  \choice \texttt{/actuator/env}
  \choice All endpoints are exposed by default
\end{choices}
\begin{solution}
  Modern Boot versions generally expose only a minimal set over HTTP by default (commonly health). Additional endpoints require explicit exposure configuration.
\end{solution}

\question[4]
When should you prefer \texttt{@SpringBootTest} over \texttt{@WebMvcTest}?
\begin{choices}
  \choice When you only want to test controller request mapping and JSON serialization in isolation
  \CorrectChoice When you need broader integration coverage across the full application context
  \choice When you want to exclude service and repository beans
  \choice Never; \texttt{@WebMvcTest} always fully replaces it
\end{choices}
\begin{solution}
  \texttt{@WebMvcTest} is a web slice; \texttt{@SpringBootTest} loads a wider context and is used for integration-style scenarios spanning multiple layers.
\end{solution}

\question[4]
Which statement about externalized configuration precedence is generally correct in Spring Boot?
\begin{choices}
  \choice Values in packaged \texttt{application.properties} always override command-line arguments
  \CorrectChoice Command-line arguments typically have high precedence and can override packaged config values
  \choice Environment variables are ignored if a properties file exists
  \choice Profile-specific files are never loaded
\end{choices}
\begin{solution}
  Boot combines multiple property sources with precedence rules; command-line options are high-priority and commonly override defaults from files.
\end{solution}

\uplevel{\par\textbf{Additional Spring Boot Concepts: Core Annotations}\par}

\question[4]
Which of the following annotations serves as the primary entry point for a Spring Boot application, bundling auto-configuration, component scanning, and extra configuration?
\begin{choices}
  \CorrectChoice \texttt{@SpringBootApplication}
  \choice \texttt{@EnableAutoConfiguration}
  \choice \texttt{@ComponentScan}
  \choice \texttt{@Configuration}
\end{choices}
\begin{solution}
  \texttt{@SpringBootApplication} is a convenience annotation that bundles \texttt{@Configuration}, \texttt{@EnableAutoConfiguration}, and \texttt{@ComponentScan}, serving as the starting point for most Spring Boot apps.
\end{solution}

\question[4]
Which annotation tells Spring Boot to automatically configure the application context based on the jar dependencies present on the classpath (e.g., setting up a DataSource if an H2 jar is found)?
\begin{choices}
  \choice \texttt{@SpringBootApplication}
  \CorrectChoice \texttt{@EnableAutoConfiguration}
  \choice \texttt{@ComponentScan}
  \choice \texttt{@Configuration}
\end{choices}
\begin{solution}
  \texttt{@EnableAutoConfiguration} drives Spring Boot's intelligent guessing mechanism to automatically configure beans that you are likely to need based on your classpath.
\end{solution}

\question[4]
Which annotation tells Spring to scan the current package and its sub-packages to find and register other Spring components (like \texttt{@Controller}, \texttt{@Service}, \texttt{@Repository})?
\begin{choices}
  \choice \texttt{@SpringBootApplication}
  \choice \texttt{@EnableAutoConfiguration}
  \CorrectChoice \texttt{@ComponentScan}
  \choice \texttt{@Configuration}
\end{choices}
\begin{solution}
  \texttt{@ComponentScan} instructs Spring to look for classes annotated with stereotypes and register them as beans in the application context.
\end{solution}

\question[4]
Which annotation marks a class as a source of bean definitions, allowing developers to manually register beans using \texttt{@Bean} methods?
\begin{choices}
  \choice \texttt{@SpringBootApplication}
  \choice \texttt{@EnableAutoConfiguration}
  \choice \texttt{@ComponentScan}
  \CorrectChoice \texttt{@Configuration}
\end{choices}
\begin{solution}
  \texttt{@Configuration} indicates that a class declares one or more \texttt{@Bean} methods and may be processed by the Spring container to generate bean definitions and service requests for those beans at runtime.
\end{solution}

\uplevel{\par\textbf{Additional Spring Boot Concepts: Data and Properties}\par}

\question[4]
How does Spring Boot natively support executing database initialization scripts (like creating tables or inserting default data) upon application startup?
\begin{choices}
  \choice By requiring the developer to write a custom \texttt{CommandLineRunner} that reads text files.
  \choice By executing any \texttt{.txt} file found in the project's root directory.
  \CorrectChoice By automatically detecting and running files named \texttt{schema.sql} and \texttt{data.sql} placed in the classpath.
  \choice Database initialization scripts are not supported and must be run manually via a database client.
\end{choices}
\begin{solution}
  Spring Boot's auto-configuration will automatically pick up \texttt{schema.sql} (for DDL) and \texttt{data.sql} (for DML) from the root of the classpath (e.g., \texttt{src/main/resources}) and execute them against the configured datasource.
\end{solution}

\question[4]
In Spring Boot, how do you typically configure a connection to an embedded or external relational database without writing boilerplate configuration classes?
\begin{choices}
  \choice By modifying the Tomcat \texttt{server.xml} configuration file.
  \CorrectChoice By providing standardized \texttt{spring.datasource.*} properties (like \texttt{url}, \texttt{username}, \texttt{password}) in the \texttt{application.properties} or \texttt{application.yml} file.
  \choice By hardcoding the database credentials directly into the repository interfaces.
  \choice Spring Boot cannot connect to a database without writing custom \texttt{@Bean} definitions for the \texttt{DataSource}.
\end{choices}
\begin{solution}
  Spring Boot's auto-configuration automatically creates a \texttt{DataSource} bean using properties prefixed with \texttt{spring.datasource} defined in the externalized configuration.
\end{solution}

\uplevel{\par\textbf{Additional practice: Spring Boot (polished Q\&A bank)}\par}

\question[4]
After you build an executable Spring Boot artifact (for example with \texttt{mvn package} or \texttt{gradle bootJar}), which command typically \textbf{runs} the fat JAR?
\begin{choices}
  \choice \texttt{java -run myapp.jar}
  \choice \texttt{maven -jar myapp.jar}
  \choice \texttt{maven run -jar myapp.jar}
  \CorrectChoice \texttt{java -jar myapp.jar}
\end{choices}
\begin{solution}
  Executable Spring Boot JARs ship a manifest \texttt{Main-Class} that delegates to \texttt{SpringApplication}; the JVM launches them with \texttt{java -jar}. Build tools produce the JAR but do not replace the runtime command shown in the bank.
\end{solution}

\question[4]
\textit{(Select all that apply.)} Which capabilities are associated with Spring Boot \textbf{DevTools}?
\begin{checkboxes}
  \CorrectChoice Automatic restart on classpath changes during development
  \choice Built-in Model--View--Controller scaffolding for every controller
  \choice A replacement IoC container separate from the main \texttt{ApplicationContext}
  \CorrectChoice Optional remote update/restart workflows for supported setups
\end{checkboxes}
\begin{solution}
  DevTools adds developer ergonomics such as restart and LiveReload integration; it is unrelated to MVC scaffolding or a second container. Remote restart support exists for specialized workflows but should be secured carefully.
\end{solution}

\endgradingrange{csaspringboot}
